% =============================================================================
% UQFF v5.78  Section Template :: T-xi
% Closure-coupling papers (xi-coupling between two or more existing closures)
% -----------------------------------------------------------------------------
% USAGE
%   1. Copy to whitepapers/PAPER_<NNNN>_<slug>.tex
%   2. Replace every <PLACEHOLDER>
%   3. Document BOTH parents (PAPER_<a> and PAPER_<b>) and the coupling xi.
% -----------------------------------------------------------------------------
% Coupling form:  xi_{ab} = <closure_a> * <weight_a> + <closure_b> * <weight_b> + ...
% Constraint:     dxi/d<param> = 0 at canonical operating point.
% =============================================================================
\documentclass[11pt]{article}
\usepackage[margin=1in]{geometry}
\usepackage{amsmath,amssymb,amsthm}
\usepackage{booktabs,longtable,array}
\usepackage{listings,xcolor}
\usepackage[colorlinks=true,linkcolor=blue,urlcolor=blue,citecolor=blue]{hyperref}
\lstset{basicstyle=\ttfamily\small,breaklines=true,frame=single,
        language=Python,showstringspaces=false,columns=fullflexible}

\title{<PAPER TITLE>: \\ $\xi$-Coupling Closure between PAPER\_<a> and PAPER\_<b>, v5.78}
\author{Star-Magic UQFF \\ PAPER\_<NNNN> \,/\, Session <SESSION>}
\date{<DATE>}

% CLOSURE :: xi_<a>_<b> :: predicted=<X> observed=<Y> error_pct=<Z>
% TEMPLATE :: T-xi
% CANONICAL :: coupling_form = xi_ab = sum_k w_k * closure_k

\begin{document}
\maketitle

\begin{abstract}
This paper derives the $\xi$-coupling between the previously closed
sectors PAPER\_<a> (\emph{<short title a>}) and PAPER\_<b>
(\emph{<short title b>}). The coupling is required to be stationary
($\partial\xi/\partial p = 0$) at the canonical operating point in
\texttt{dpm\_vacuum\_manifold.py} v3.0, and we verify the resulting
joint prediction against <observational anchor>.
\end{abstract}

\section{Parent Closures}
\begin{description}
  \item[PAPER\_<a>] <short title>, closure $C_a = <value>$, $z_a=<z>$.
  \item[PAPER\_<b>] <short title>, closure $C_b = <value>$, $z_b=<z>$.
\end{description}

\section{Coupling Definition}
\begin{equation}
  \xi_{ab} \;=\; w_a\, C_a \;+\; w_b\, C_b \;+\; \kappa_{ab}\, C_a C_b
\end{equation}
with weights $w_a, w_b$ derived from the canonical $|SO(5)|=10$ ratio and
$\kappa_{ab}$ proportional to $F_{TRZ}^{\,n}$ where $n=<n>$ is fixed by
dimensional analysis.

\section{Stationarity Constraint}
At the canonical operating point,
\begin{equation}
  \left.\frac{\partial \xi_{ab}}{\partial p}\right|_{\,p=p_\star} \;=\; 0
\end{equation}
which fixes <which canonical parameter or relationship>.

\section{Joint Prediction}
\begin{longtable}{lll}
\toprule
Observable & Predicted (joint) & Observed (cite) \\
\midrule
<obs> & <value> & <value $\pm$ error> \\
\bottomrule
\end{longtable}
$\varepsilon = <Z>\%$, $z = <z>$.

\section{Cross-references}
\begin{itemize}
  \item Parent PAPER\_<a>, parent PAPER\_<b>.
  \item Ledger row: \texttt{xi\_<a>\_<b>} in \texttt{master\_closures.csv}.
\end{itemize}

\bibliographystyle{plain}
\begin{thebibliography}{9}
\bibitem{ref1} <DOI/arXiv for the joint observational anchor>.
\bibitem{paperA} PAPER\_<a> (Star-Magic).
\bibitem{paperB} PAPER\_<b> (Star-Magic).
\end{thebibliography}

\end{document}
